\section{Introduction}
\label{sec:intro}

Peatland and forest fires in Indonesia produce some of the densest smoke on Earth. During strong
El Niño dry seasons, haze from Sumatra and Kalimantan covers large parts of Equatorial Southeast
Asia for weeks, closing schools and airports and causing severe respiratory illness
\citep{huang2024haze}. Communities living around forests and peatland often learn that smoke has
arrived only when it is already outside their window. Short-term haze forecasts that are local,
open and honest about their uncertainty could help these communities prepare.

Satellite fire detections are the obvious predictor. NASA's Fire Information for Resource
Management System (FIRMS) distributes VIIRS active-fire detections within hours of each overpass
\citep{firms}. Smoke, however, travels with the wind. A detection 80\,km upwind of a town matters
more than one 80\,km downwind. Wind-weighted fire-exposure features are therefore an intuitive and
increasingly common design choice \citep{gnnfire2023,sci2026exposure}, yet their \emph{directional}
component is rarely tested on its own. When a model with wind-weighted fire exposure beats a
model without fire information, the gain may come from the fire amount, distance or recency rather
than from where the fires are relative to the wind. Other studies find that fire predictors add
little once recent pollution and weather are known \citep{santacruz2026}.

A second obstacle is the target. Hourly PM\textsubscript{2.5} observations over Kalimantan and most
of Sumatra are short, sparse or unavailable before 2021. We instead use a long and openly
archived proxy: \emph{airport visibility}. Routine aviation weather reports (METAR) record
horizontal visibility and present-weather codes for haze (HZ) and smoke (FU) every 30--60 minutes
and have been archived for decades \citep{iem}. Visibility below 3\,km with a haze or smoke code
and no precipitation or fog is a direct, observed marker of smoke-haze impact.

We ask one question with a matched design. \textbf{Do wind-directional satellite fire features
improve 6\,h and 24\,h haze-visibility forecasts beyond a non-directional fire-exposure feature with
identical distance, age and fire-radiative-power weighting?} We also ask whether any gain survives
realistic fire-data latency. Our contributions:

\begin{enumerate}
  \item An eleven-year (2015--2025), twelve-airport benchmark of haze-limited visibility in Borneo
        and Sumatra, built entirely from keyless public archives, covering three major haze
        seasons (2015, 2019, 2023).
  \item A matched ablation that isolates the directional term (directional vs non-directional
        exposure) against a strong station-state baseline that already includes the local wind
        direction, with time-shifted and mask-preserving permuted-wind placebos, missingness and
        wind-report-category controls. We show why rotated-wind placebos are not clean nulls.
  \item Leave-one-year-out evaluation with station- and region-level block bootstrap, breakdowns by
        season and reporting convention, a vegetation-only sensitivity and a fire-data latency
        analysis, and evidence that single fits of tree models cannot resolve effects of this
        size in this setting.
  \item An open implementation used in a public smoke-awareness application, Pavana, whose
        deployed model is the feature set evaluated here.
\end{enumerate}

The answer is negative, and reaching it taught us two lessons about method. In a pre-specified
comparison of single gradient-boosted fits, adding the directional exposure raised 24\,h average
precision from \pv{loyo24.AP.D0} to \pv{loyo24.AP.E}, a large and significant gain. Placebo winds
reproduced almost all of it, and so did refitting the same model with another random seed: the
differences between single fits were dominated by seed-to-seed variation in one extreme season.
With five-seed ensembles, the directional model is at most a few thousandths of average precision
above time-shifted and permuted-wind placebos and non-directional controls with the same
missingness, with no consistent sign across aggregations, horizons or the two model families. Fire counts improve 24\,h forecasts modestly. We think
the checks we needed, seed ensembles and mask-preserving placebos, apply to any study that credits
wind-weighted fire features.

The protocol, including the primary hypothesis, was written down and committed before any result
reported here was computed (Appendix~\ref{app:protocol}). Changes made after seeing results are
listed there and labelled in the text.
